\chapter{Solutions to Selected Exercises}%
\label{app:solutions}

Here are worked solutions to a selection of exercises from across the book. They
are not the only correct answers programming rewards many paths to the same
result but each is complete, idiomatic, and runs as shown. Try the exercise
yourself before reading the solution; the struggle is where the learning happens.

\section*{Chapter 1 Getting Started}

\textbf{Exercise 1.2} (print a labelled calculation):

\begin{salamen}
func main:
    println "12 * 8 =", 12 * 8
end
\end{salamen}

\begin{progout}
12 * 8 = 96
\end{progout}

\section*{Chapter 2 Variables and Data Types}

\textbf{Exercise 2.3} (average as a float):

\begin{salamen}
func main:
    sum : int = 10 + 15 + 20
    avg : f64 = (sum as f64) / 3.0
    println "average =", avg
end
\end{salamen}

\begin{progout}
average = 15
\end{progout}

The casts to \code{f64} are what make this a real division; without them
\code{45 / 3} would (here, coincidentally) also give 15, but for non-divisible
totals integer division would be wrong.

\section*{Chapter 3 Operators}

\textbf{Exercise 3.3} (between 1 and 12 inclusive):

\begin{salamen}
func main:
    n : int = 7
    in_range : auto = n >= 1 && n <= 12
    println n, "in 1..12?", in_range
end
\end{salamen}

\begin{progout}
7 in 1..12? true
\end{progout}

\textbf{Exercise 3.6} (cube in place with \code{\textasciicircum\textasciicircum=}):

\begin{salamen}
func main:
    mut side : f64 = 3.0
    side ^^= 3
    println side
end
\end{salamen}

\begin{progout}
27
\end{progout}

Starting from \code{mut side : int = 3} instead gives a compile error:
\code{\textasciicircum\textasciicircum} always produces a \code{f64}, and an \code{f64} cannot be assigned
back into an \code{int} variable without an explicit \code{as int} cast, so
\code{side \textasciicircum\textasciicircum= 3} is rejected the same way \code{side = side \textasciicircum\textasciicircum 3} would be.

\section*{Chapter 4 Control Flow}

\textbf{Exercise 4.3} (the sign of a number):

\begin{salamen}
func sign(n : int) : str:
    if n < 0: ret "negative"
    else n == 0: ret "zero"
    else: ret "positive"
    end
end

func main:
    println sign(-5), sign(0), sign(42)
end
\end{salamen}

\begin{progout}
negative zero positive
\end{progout}

\section*{Chapter 5 Loops}

\textbf{Exercise 5.2} (sum 1 to 100):

\begin{salamen}
func main:
    mut total : int = 0
    repeat 1 to 100 with i:
        total += i
    end
    println total
end
\end{salamen}

\begin{progout}
5050
\end{progout}

\section*{Chapter 6 Functions}

\textbf{Exercise 6.4} (recursive sum):

\begin{salamen}
func sum_to(n : int) : int:
    if n <= 0: ret 0 end       // base case
    ret n + sum_to(n - 1)      // recursive step
end

func main:
    println sum_to(10)
end
\end{salamen}

\begin{progout}
55
\end{progout}

\section*{Chapter 7 Arrays and Vectors}

\textbf{Exercise 7.3} (a vector of squares):

\begin{salamen}
func main:
    mut squares : Vector<int> = Vector {}
    repeat 1 to 10 with i:
        squares.push(i * i)
    end
    println "count:", squares.len()
    mut j : int = 0
    until j < squares.len():
        print squares.get(j)[0]
        print " "
        j += 1
    end
    println ""
    squares.free()
end
\end{salamen}

\begin{progout}
count: 10
1 4 9 16 25 36 49 64 81 100
\end{progout}

\section*{Chapter 8 Strings}

\textbf{Exercise 8.4} (palindrome test):

\begin{salamen}
import "str"

func is_palindrome(s : str) : bool:
    ret str.Equals(s, str.Reverse(s))
end

func main:
    println is_palindrome("level")
    println is_palindrome("salam")
end
\end{salamen}

\begin{progout}
true
false
\end{progout}

\section*{Chapter 9 Structs and Methods}

\textbf{Exercise 9.1} (a circle's area):

\begin{salamen}
struct Circle:
    radius : f64 = 0.0
    func area() : f64:
        ret 3.14159 * this.radius * this.radius
    end
end

func main:
    c : Circle = Circle { radius = 2.0 }
    println "area:", c.area()
end
\end{salamen}

\begin{progout}
area: 12.5664
\end{progout}

\section*{Chapter 10 Enums}

\textbf{Exercise 10.3} (turn right):

\begin{salamen}
enum Direction: North, East, South, West end

func turn_right(d : Direction) : Direction:
    if d == Direction.North: ret Direction.East
    else d == Direction.East: ret Direction.South
    else d == Direction.South: ret Direction.West
    else: ret Direction.North
    end
end

func main:
    mut d : Direction = Direction.North
    d = turn_right(d)
    println "now facing index:", d as int      // East = 1
end
\end{salamen}

\begin{progout}
now facing index: 1
\end{progout}

\section*{Chapter 11 Interfaces}

\textbf{Exercise 11.1} (a greeter interface):

\begin{salamen}
interface Greeter:
    func greet() : str
end

struct Formal:
    func greet() : str: ret "Good evening." end
end

struct Casual:
    func greet() : str: ret "Hey!" end
end

func announce<T: Greeter>(g : T):
    println g.greet()
end

func main:
    announce(Formal {})
    announce(Casual {})
end
\end{salamen}

\begin{progout}
Good evening.
Hey!
\end{progout}

\section*{Chapter 12 Generics}

\textbf{Exercise 12.1} (a generic \code{Min}):

\begin{salamen}
func Min<T>(a : T, b : T) : T:
    if a < b: ret a end
    ret b
end

func main:
    println Min(10, 20)
    println Min(2.5, 1.5)
end
\end{salamen}

\begin{progout}
10
1.5
\end{progout}

\section*{Chapter 13 Closures}

\textbf{Exercise 13.3} (a multiplier factory):

\begin{salamen}
func make_multiplier(factor : int) : func(int) int:
    ret (x : int) => x * factor
end

func main:
    times3 : auto = make_multiplier(3)
    times10 : auto = make_multiplier(10)
    println times3(5)
    println times10(5)
end
\end{salamen}

\begin{progout}
15
50
\end{progout}

\section*{Chapter 14 Errors and Defer}

\textbf{Exercise 14.2} (find the first negative):

\begin{salamen}
func find_first_negative(a : int[5]) : int:
    repeat 5 with i:
        if a[i] < 0:
            ret i
        end
    end
    ret -1
end

func main:
    data : int[5] = [3, 8, -2, 5, -7]
    println "first negative at:", find_first_negative(data)
end
\end{salamen}

\begin{progout}
first negative at: 2
\end{progout}

\section*{Chapter 16 Standard Library, Part 1}

\textbf{Exercise 16.3} (the hypotenuse):

\begin{salamen}
import "math"

func main:
    a : f64 = 3.0
    b : f64 = 4.0
    println "hypotenuse:", math.Sqrt(a * a + b * b)
end
\end{salamen}

\begin{progout}
hypotenuse: 5
\end{progout}

\section*{Chapter 17 Standard Library, Part 2}

\textbf{Exercise 17.1} (a map of ages):

\begin{salamen}
func main:
    mut ages : HashMap<str, int> = HashMap {}
    ages.insert("Sara", 30)
    ages.insert("Ali", 25)
    ages.insert("Mina", 41)
    println "Ali is", ages.get("Ali")
    println "have Reza?", ages.has("Reza")
    ages.free()
end
\end{salamen}

\begin{progout}
Ali is 25
have Reza? false
\end{progout}

\section*{Chapter 21 Testing}

\textbf{Exercise 21.1} (testing factorial):

\begin{salamen}
import "testing"

func factorial(n : int) : int:
    if n <= 1: ret 1 end
    ret n * factorial(n - 1)
end

func test_factorial():
    testing.AssertEqInt(factorial(0), 1)
    testing.AssertEqInt(factorial(5), 120)
    testing.AssertEqInt(factorial(3), 6)
end

func main:
    test_factorial()
    testing.Summary()
end
\end{salamen}

\begin{progout}
tests: 3 passed, 0 failed
\end{progout}

\section*{Beyond these}

The remaining exercises follow the same patterns shown here and in the chapters.
If you get stuck, re-read the relevant section, run a smaller version of the
problem, and add \code{println} probes to see what your code is actually doing ---
the very habits from Chapter~21. Every program in this book was written by someone
who started exactly where you are now.
